\SourceLeaf{586}\pdfbookmark[1]{Astronomische Nachrichten No. 417. 1841. February 9.}{nachrichten-417}
\begin{center}\small\textsc{Schumacher}.\quad Astronomische Nachrichten No.\ 417.\quad 1841.\quad February 9.\end{center}

I recently read in a newspaper that the American naval captain \textsc{Wilkes} had come fairly close to
\begin{center}\textit{the magnetic south pole}\end{center}
and that at $67^0\;4'$ south latitude and $147^0\;30'$ longitude (doubtless east of Greenwich) he had found the magnetic declination to be $12^0\;35'$ east and the inclination $87^0\;30'$. On this basis, and after a quick calculation, I would provisionally place the actual pole
\begin{center}at $70^0\;21'$ south latitude,\qquad $146^0\;17'$ longitude.\end{center}
This position is much nearer the one my theory gave [p.\ 163 of this volume] than I myself had expected. The actual pole, as I had supposed there, farther north than the one calculated from the theory. But the difference in latitude amounts to only one third of the difference I was prepared for after looking at the observations from Hobarttown. The actual pole likewise lies farther west than the one calculated from the theory, and here the difference is almost exactly the size I had presumed in that passage. It is unnecessary to remark that at these high latitudes a difference of six degrees of longitude amounts to only as much as two degrees of latitude.

